\section{The empirical distribution function}

The cumulative distribution function of a random variable is
\[
F_X(x)=P(X\leq x).
\]
It tells us how much model probability lies to the left of \(x\). For a dataset,
the analogous object is the empirical distribution function. It tells us what
proportion of the observed data lies to the left of \(x\).

\begin{definitionbox}
For a dataset \(x_1,\ldots,x_n\), the \textbf{empirical distribution function},
or \textbf{empirical CDF}, is
\[
F_n(x)=\frac1n\sum_{i=1}^n\mathbf{1}_{\{x_i\leq x\}}.
\]
\end{definitionbox}

The value \(F_n(x)\) is the proportion of observations not larger than \(x\). It
is the CDF of the empirical distribution.

\subsection*{A hand calculation}

Consider the dataset
\[
2,\quad 4,\quad 4,\quad 7,\quad 9.
\]
Here \(n=5\). Let us compute several values of \(F_n\).

For \(x=3\), only the observation \(2\) is at most \(3\), so
\[
F_n(3)=\frac15.
\]
For \(x=4\), the observations \(2,4,4\) are at most \(4\), so
\[
F_n(4)=\frac35.
\]
For \(x=8\), the observations \(2,4,4,7\) are at most \(8\), so
\[
F_n(8)=\frac45.
\]
For \(x=10\), all observations are at most \(10\), so
\[
F_n(10)=1.
\]

\subsection*{Step function interpretation}

The empirical CDF is a step function. It is constant between observed values and
jumps at observed values. In the example above, it jumps at \(2\), \(4\), \(7\),
and \(9\). At the value \(4\), it jumps by \(2/5\), because two observations are
equal to \(4\).

This is the graphical form of the empirical distribution. Each observation
contributes mass \(1/n\), and the empirical CDF accumulates that mass from left
to right.

\begin{remarkbox}
The empirical CDF is always a step function for a finite dataset. The theoretical
CDF of a continuous model may be smooth, but the empirical CDF of a finite sample
from that model is still a step function.
\end{remarkbox}

\subsection*{Comparing \(F_n\) with \(F_X\)}

If the data are generated from a theoretical distribution with CDF \(F_X\), then
we can compare the empirical CDF \(F_n\) with \(F_X\). The two functions answer
related but different questions:
\[
F_X(x)=P(X\leq x)
\]
is a model probability, while
\[
F_n(x)=\frac{\#\{i:x_i\leq x\}}{n}
\]
is an observed proportion.

For example, if the theoretical model is uniform on \([0,1]\), then
\[
F_X(x)=x,
\qquad 0\leq x\leq1.
\]
A sample of size \(20\) from this model gives an empirical CDF with jumps of size
\(1/20\). It will not usually equal the line \(F_X(x)=x\), but it may be close in
many places. A larger sample often gives a closer visual approximation.

This observation is important, but at this stage it remains descriptive. Later,
limit theorems will explain why empirical summaries stabilize under suitable
assumptions.

\subsection*{Interval proportions from the empirical CDF}

Just as a theoretical CDF gives interval probabilities by differences, the
empirical CDF gives interval proportions. For \(a<b\),
\[
F_n(b)-F_n(a)
\]
is the proportion of observations satisfying
\[
a<x_i\leq b.
\]
Thus
\[
F_n(b)-F_n(a)=\frac{\#\{i:a<x_i\leq b\}}{n}.
\]
Endpoint conventions matter for empirical data because observations can be
exactly equal to endpoints. This is different from continuous density models,
where individual endpoints usually have probability zero.

\begin{summarybox}
The empirical CDF is
\[
F_n(x)=\frac1n\sum_{i=1}^n\mathbf{1}_{\{x_i\leq x\}}.
\]
It is the observed proportion of data not exceeding \(x\). It is the data-based
counterpart of the theoretical CDF \(F_X(x)=P(X\leq x)\).
\end{summarybox}
